% GENERATED FILE. Edit canonical lessons, then run npm run generate:books.
% canonical-source-hash: fnv1a64:172e2f295f0cb32f
% canonical-lessons: PT-C07-dias-1, PT-C07-dias-2

\chapter{Days of the Week}
\label{ch:pt-days-of-the-week}
\hlchaptermodality{7}

\begin{chapteropening}
\textbf{By the end of this chapter:} \emph{I can name all seven Portuguese days and count the weekdays from domingo the way Portuguese does.}

Run domingo, segunda, terça, quarta, quinta, sexta, sábado, and explain why Monday is the second day: domingo, the Lord's day, is counted first.
\end{chapteropening}

\section[segunda, terça, quarta, quinta …]{segunda to sexta — the week that got numbered}

\label{lesson:PT-C07-dias-1}



\begin{quote}
Here Portuguese does something \textbf{no other Romance language does}. French, Italian, and Spanish all name their weekdays for planet-gods (\emph{lundi}, \emph{lunedì}, \emph{lunes} = the Moon). Portuguese threw the pagan gods out and \textbf{numbered} the days instead — and the numbers are the very ones you just learned.
\end{quote}

\begin{culture}
In the 6th century, \textbf{Martinho de Braga}, a bishop in what is now northern Portugal, refused to let Christians name days after \emph{Mars} and \emph{Jupiter}. Borrowing the Church's word for the numbered holy days of Easter week, he called them \textbf{feiras} — and Portuguese alone kept the system for \textbf{every} week.

\textbf{feira} $\leftarrow$ Latin \textbf{fēria}, \textquotedblleft{}a feast-day / a day free of work.\textquotedblright{} (The same \emph{fēria} drifted, elsewhere, into \textquotedblleft{}market/fair\textquotedblright{} — English \textbf{fair} is its cousin, because markets were held on feast-days.)
\end{culture}

\begin{grammarlens}[title={Grammar Lens: The pattern: {[}ordinal{]} + feira}]
\noindent
\begin{tabularx}{\linewidth}{@{}>{\raggedright\arraybackslash}X>{\raggedright\arraybackslash}X>{\raggedright\arraybackslash}X>{\raggedright\arraybackslash}X@{}}
\toprule
\textbf{Portuguese} & \textbf{means} & \textbf{$\leftarrow$ ordinal of} & \textbf{(cardinal you learned)} \\
\midrule
\textbf{segunda}-feira (Mon) & \textquotedblleft{}2nd feast-day\textquotedblright{} & \emph{secundus} \textquotedblleft{}second\textquotedblright{} & \emph{dois} (2) \\
\textbf{terça}-feira (Tue) & \textquotedblleft{}3rd\textquotedblright{} & \emph{tertius} \textquotedblleft{}third\textquotedblright{} & \emph{três} (3) \\
\textbf{quarta}-feira (Wed) & \textquotedblleft{}4th\textquotedblright{} & \emph{quartus} \textquotedblleft{}fourth\textquotedblright{} & \emph{quatro} (4) \\
\textbf{quinta}-feira (Thu) & \textquotedblleft{}5th\textquotedblright{} & \emph{quintus} \textquotedblleft{}fifth\textquotedblright{} & \emph{cinco} (5) \\
\textbf{sexta}-feira (Fri) & \textquotedblleft{}6th\textquotedblright{} & \emph{sextus} \textquotedblleft{}sixth\textquotedblright{} & \emph{seis} (6) \\
\bottomrule
\end{tabularx}

So \emph{quarta-feira} is literally \textquotedblleft{}\textbf{fourth} feast-day,\textquotedblright{} and you can hear \emph{quatro} (4) inside \emph{quarta}. In everyday speech the \textquotedblleft{}-feira\textquotedblright{} is often dropped: \textquotedblleft{}\emph{Até segunda!}\textquotedblright{} — \textquotedblleft{}See you Monday!\textquotedblright{} (Why does Monday start at \textbf{two}, not one? Because Sunday is counted as the \emph{first} day — the next lesson closes that loop.)
\end{grammarlens}

\subsection*{Your turn}
Say these aloud:

\begin{itemize}
\raggedright
  \item \textquotedblleft{}segunda, terça, quarta, quinta, sexta\textquotedblright{} — Monday through Friday
  \item hear the number in each — \textquotedblleft{}quarta $\leftarrow$ quatro (4), quinta $\leftarrow$ cinco (5), sexta $\leftarrow$ seis (6)\textquotedblright{}
  \item drop the tail — \textquotedblleft{}Até sexta!\textquotedblright{} (\textquotedblleft{}See you Friday!\textquotedblright{})
\end{itemize}

\begin{tcolorbox}[breakable,colback=teal!4,colframe=teal!35!black,title={Before you move on}]
How is Portuguese unique among Romance languages here? (It \textbf{numbers} its weekdays — \emph{feiras} — instead of naming planet-gods.) What does \emph{feira} mean, and its English cousin? (\textquotedblleft{}Feast-day,\textquotedblright{} $\leftarrow$ \emph{fēria}; English \emph{fair}.) What number hides in \emph{quinta-feira}? (Five — \emph{quinta} $\leftarrow$ \emph{quintus}, cousin of \emph{cinco}.) Next: the two days that \textbf{kept} their old names — \emph{sábado} and \emph{domingo}.
\end{tcolorbox}
\section[sábado, domingo]{sábado and domingo — and why Monday is \textquotedblleft{}the second\textquotedblright{}}

\label{lesson:PT-C07-dias-2}



\begin{quote}
Five weekdays got numbered — but two kept their names. \emph{Sábado} and \emph{domingo} are the \textbf{religious anchors} of the Portuguese week, and \emph{domingo} is the key that explains why the counting started where it did.
\end{quote}

\begin{cousinweb}
\noindent
\begin{tabularx}{\linewidth}{@{}>{\raggedright\arraybackslash}X>{\raggedright\arraybackslash}X>{\raggedright\arraybackslash}X>{\raggedright\arraybackslash}X@{}}
\toprule
\textbf{Portuguese} & \textbf{$\leftarrow$ Latin} & \textbf{meaning} & \textbf{shared with} \\
\midrule
\textbf{sábado} & \emph{Sabbatum} $\leftarrow$ Hebrew \emph{shabbāt} & the \textquotedblleft{}\textbf{Sabbath},\textquotedblright{} rest-day & Spanish \emph{sábado}, Italian \emph{sabato}, French \emph{samedi} \\
\textbf{domingo} & \emph{(diēs) Dominica} $\leftarrow$ \emph{Dominus} & \textquotedblleft{}the \textbf{Lord's} day\textquotedblright{} & Spanish \emph{domingo}, Italian \emph{domenica} \\
\bottomrule
\end{tabularx}

\begin{itemize}
\raggedright
  \item \textbf{sábado} is the \textbf{Sabbath}, from Hebrew \emph{shabbāt} \textquotedblleft{}to rest\textquotedblright{} — the one weekend name every Romance language shares (even numbering-mad Portuguese kept it).
  \item \textbf{domingo} is \textquotedblleft{}the \textbf{Lord's} day,\textquotedblright{} \emph{diēs Dominica}, from \emph{Dominus} \textquotedblleft{}lord, master\textquotedblright{} ($\to$ English \emph{dominion, dominate, dame}).
\end{itemize}
\end{cousinweb}

\begin{grammarlens}[title={Grammar Lens: domingo is \textquotedblleft{}day one\textquotedblright{} — so Monday is \textquotedblleft{}day two\textquotedblright{}}]
Now the loop closes. The Church counted \textbf{domingo (Sunday) as the \emph{first} day} of the week — \emph{prima feria}, the Lord's day, the start. That's why:

\begin{quote}
domingo = 1st $\to$ \textbf{segunda}-feira (Monday) = \textbf{2nd} $\to$ terça (3rd) $\to$ … $\to$ sexta (Friday) = 6th $\to$ sábado (Saturday) = 7th.
\end{quote}

Monday isn't the first weekday in this system — it's the day \textbf{after} the Lord's day, the \emph{second}. So \emph{segunda-feira} (\textquotedblleft{}second feast-day\textquotedblright{}) makes perfect sense once you know Sunday was counted first. The whole week is a countdown anchored on \emph{domingo}.
\end{grammarlens}

\subsection*{Your turn}
Say these aloud:

\begin{itemize}
\raggedright
  \item the full week — \textquotedblleft{}domingo, segunda, terça, quarta, quinta, sexta, sábado\textquotedblright{}
  \item the anchors — \textquotedblleft{}sábado = Sabbath, domingo = the Lord's day\textquotedblright{}
  \item the logic — \textquotedblleft{}domingo is 1st, so segunda (Monday) is 2nd\textquotedblright{}
\end{itemize}

\begin{tcolorbox}[breakable,colback=teal!4,colframe=teal!35!black,title={Before you move on}]
Give all seven Portuguese days. (\emph{Domingo, segunda … sábado}.) What does \emph{domingo} mean, and why does that make Monday \textquotedblleft{}\emph{segunda}\textquotedblright{}? (\textquotedblleft{}The \textbf{Lord's} day\textquotedblright{}; Sunday is counted \textbf{first}, so Monday is the \textbf{second}.) Which weekend name do all the Romance languages share? (\emph{Sábado} — the Sabbath.) Next chapter builds on this toward telling the time.
\end{tcolorbox}
